\documentclass[11pt]{article}

\usepackage{amsmath,amssymb,amsthm}
\usepackage{mathtools}
\usepackage[margin=1in]{geometry}

\newtheorem{theorem}{Theorem}
\newtheorem{lemma}{Lemma}
\newtheorem{proposition}{Proposition}
\newtheorem{corollary}{Corollary}
\newtheorem{definition}{Definition}
\newtheorem{remark}{Remark}

\DeclareMathOperator{\vN}{vN}
\DeclareMathOperator{\Tr}{Tr}
\DeclareMathOperator{\Ad}{Ad}
\DeclareMathOperator{\link}{lk}
\DeclareMathOperator{\str}{star}
\DeclareMathOperator{\WOT}{WOT}
\newcommand{\Gra}{\Gamma}
\newcommand{\VV}{\mathcal V}
\newcommand{\EE}{\mathcal E}
\newcommand{\BB}{\mathbb B}
\newcommand{\KK}{\mathbb K}
\newcommand{\HH}{\mathcal H}
\newcommand{\Lt}{L^2(M)}
\newcommand{\veps}{\varepsilon}
\newcommand{\wstar}{\mathrm{w}^*}
\newcommand{\ip}[1]{\langle #1\rangle}
\newcommand{\norm}[1]{\lVert #1\rVert}
\newcommand{\abs}[1]{\lvert #1\rvert}
\newcommand{\set}[1]{\{#1\}}
\newcommand{\rest}{\!\restriction\!}

\begin{document}

\begin{center}
{\Large\bfseries Proper proximality of irreducible graph product von Neumann algebras}
\end{center}

\noindent\textbf{Status: partial (rigorous reduction + base case complete; one permanence step cited).}
The solution gives complete, self-contained proofs of: (i) the join/irreducibility dichotomy and its
operator-algebraic meaning (a join is exactly a tensor splitting, which is the obstruction to proper
proximality); (ii) the base case of edgeless graphs (free products), where proper proximality is
established from the Ding--Elayavalli--Peterson free-product theorem; and (iii) the reduction of the
general irreducible case to a single amalgamated--free--product permanence statement
(Theorem~\ref{thm:perm-afp}) via the link decomposition $M=M_{\str(v)}*_{M_{\link(v)}}M_{\Gra\setminus v}$,
together with a verification, using the trace-zero unitaries and irreducibility, that the malnormality /
weak-mixing hypotheses of that permanence statement hold. The one ingredient that is invoked rather than
reproved from first principles is the amalgamated--free--product permanence theorem for proper proximality
(Theorem~\ref{thm:perm-afp}); its exact hypotheses are stated precisely and each is verified in our setting.
The conditional logic ``Theorem~\ref{thm:perm-afp} $\Rightarrow$ main result'' is complete and gap-free.

\bigskip

\section{Statement and conventions}

Throughout, a \emph{tracial von Neumann algebra} is a pair $(N,\tau)$ with $N$ a von Neumann algebra on a
separable Hilbert space and $\tau$ a faithful normal tracial state. We write $L^2(N)=L^2(N,\tau)$ for the
GNS Hilbert space, $\hat x\in L^2(N)$ for the image of $x\in N$, and $\hat 1$ for the cyclic trace vector.
A unitary $u\in N$ is \emph{trace-zero} if $\tau(u)=0$.

Let $\Gra=(\VV,\EE)$ be a finite simplicial graph: $\VV=\VV(\Gra)$ is a finite vertex set and
$\EE=\EE(\Gra)\subseteq\binom{\VV}{2}$ is a set of unordered edges; no loops, no multiple edges. For
$u,v\in\VV$ we write $u\sim v$ if $\{u,v\}\in\EE$ and $u\not\sim v$ otherwise (with $v\not\sim v$). For
$v\in\VV$ put
\[
\link(v)=\set{w\in\VV : w\sim v},\qquad \str(v)=\link(v)\cup\set v .
\]
For $S\subseteq\VV$ let $\Gra[S]$ be the induced subgraph on $S$.

\begin{definition}[Join and irreducibility]\label{def:join}
$\Gra$ is a \emph{(nontrivial) join} if there are nonempty $\VV_1,\VV_2$ with
$\VV_1\cup\VV_2=\VV$, $\VV_1\cap\VV_2=\emptyset$, and $v\sim u$ for all $v\in\VV_1$, $u\in\VV_2$;
we write $\Gra=\Gra[\VV_1]+\Gra[\VV_2]$. We say $\Gra$ is \emph{irreducible} if it is not a nontrivial join.
\end{definition}

\begin{remark}\label{rem:join-statement}
The problem's phrasing of a join asks for subgraphs $\Gra_1,\Gra_2$ with
$\VV(\Gra_1)\cup\VV(\Gra_2)=\VV$ and $v\sim u$ for all $v\in\VV(\Gra_1),u\in\VV(\Gra_2)$. If
$\VV(\Gra_1)\cap\VV(\Gra_2)\ne\emptyset$, choosing $w$ in the intersection forces $w\sim w$, impossible in a
simplicial graph; so any genuine join has disjoint vertex parts, matching Definition~\ref{def:join}. A join
is required nontrivial (both parts nonempty), as in the statement.
\end{remark}

\subsection*{Graph products}

For tracial $(M_v,\varphi_v)_{v\in\VV}$ (here $\varphi_v=\tau_v$ is the trace), the
\emph{graph product} $(M,\tau)=(\ast_{\Gra,v}M_v,\ \ast_{\Gra}\tau_v)$ is the tracial von Neumann algebra
generated by copies of the $M_v$ subject to: copies $M_v,M_w$ \emph{commute} whenever $v\sim w$, and are
\emph{freely independent} (with respect to the trace) whenever $v\not\sim w$. We use the standard model and
the word-length orthogonal decomposition of $L^2(M)$, recalled in Section~\ref{sec:structure}. We write
$M_S=\vN(M_v:v\in S)\subseteq M$ for $S\subseteq\VV$; this is canonically the graph product over the induced
subgraph $\Gra[S]$, and there is a trace-preserving conditional expectation $E_S:M\to M_S$.

\medskip

\begin{theorem}[Main result]\label{thm:main}
Let $\Gra$ be a finite simplicial graph with $\abs{\VV}\ge 3$, let $(M_v,\tau_v)$ be separable tracial von
Neumann algebras, and suppose:
\begin{enumerate}
\item[(a)] $\Gra$ is irreducible (Definition~\ref{def:join}); and
\item[(b)] each $M_v$ contains a trace-zero unitary.
\end{enumerate}
Then the graph product $M=\ast_{\Gra,v}M_v$ is properly proximal.
\end{theorem}

\section{Proper proximality for von Neumann algebras}\label{sec:pp}

We use the framework of Ding--Elayavalli--Peterson \cite{DEP}, which extends Boutonnet--Ioana--Peterson
\cite{BIP} from groups to tracial von Neumann algebras. We recall exactly what is needed.

Let $(M,\tau)$ be tracial, acting standardly on $\Lt$ with the left action; let $J$ be the modular
conjugation and $M^{\mathrm{op}}=JMJ$ the commutant (the right action). Let $\KK=\KK(\Lt)$ be the compact
operators. Consider the unital $C^*$-subalgebra of $\BB(\Lt)$
\begin{equation}\label{eq:smallinfty}
\mathbb S(M)=\set{T\in\BB(\Lt) : [T,\,J x J]\in\KK\ \text{for all } x\in M }
\end{equation}
of operators that asymptotically commute with the right action; equivalently $T$ is \emph{small at
infinity} from the right. This is a $C^*$-algebra containing $\KK$ and is an $M$--$M$ operator bimodule via
$(a,b)\cdot T=aTb$ ($a,b\in M$, where $M$ acts on the left and $M^{\mathrm{op}}=JMJ$ acts as the right
$M$-action). Set $\mathbb S(M)_0=\mathbb S(M)\cap (\mathbb S(M))^{*}$ and let
\begin{equation}\label{eq:quotient}
\widetilde{\mathbb S}(M)=\mathbb S(M)/\KK,
\qquad q:\mathbb S(M)\to\widetilde{\mathbb S}(M)\ \text{the quotient.}
\end{equation}
A state $\theta$ on $\mathbb S(M)$ is \emph{left-$M$-central} (resp.\ \emph{right-$M$-central}) if
$\theta(aT)=\theta(Ta)$ for all $a\in M$ (resp.\ $\theta((JaJ)T)=\theta(T(JaJ))$). A state is
\emph{boundary} if it vanishes on $\KK$ (factors through $q$).

\begin{definition}[Proper proximality, \cite{DEP}]\label{def:pp}
A tracial von Neumann algebra $(M,\tau)$ is \emph{not properly proximal} if there exists a state $\theta$ on
$\mathbb S(M)$ which
\begin{enumerate}
\item[(i)] is \emph{normal} as a left $M$-functional and as a right $M$-functional, with
$\theta\rest_{M}=\tau=\theta\rest_{M^{\mathrm{op}}}$ (it restricts to the vector state of $\hat 1$ on each
copy of $M$); and
\item[(ii)] is left-$M$-central;
\end{enumerate}
\emph{and} which is moreover a \emph{boundary} state, i.e.\ $\theta\rest_{\KK}=0$. Equivalently, $M$ is
\emph{properly proximal} if there is \textbf{no} left-$M$-central boundary state on
$\widetilde{\mathbb S}(M)$ that restricts to $\tau$ on both copies of $M$.
\end{definition}

Definition~\ref{def:pp} is the $C^*$-bimodule reformulation of \cite[\S3]{DEP}. We record the three
permanence facts we use; each is a theorem of \cite{DEP} (the group analogues are in \cite{BIP}). We state
them precisely and will verify their hypotheses in our setting.

\begin{theorem}[Tensor obstruction \cite{DEP}]\label{thm:tensor-obstruct}
If $M=M_1\bar\otimes M_2$ with both $M_1,M_2$ diffuse, then $M$ is \emph{not} properly proximal.
\end{theorem}

\begin{theorem}[Free products \cite{DEP}]\label{thm:freeprod}
Let $(N_1,\tau_1),\dots,(N_k,\tau_k)$ ($k\ge2$) be tracial von Neumann algebras, each containing a
trace-zero unitary (in particular each diffuse), and let $N=N_1*\cdots*N_k$ be their tracial free product.
Then $N$ is properly proximal, and moreover \emph{properly proximal relative to $\mathbb C$} (i.e.\ a boundary
state as in Definition~\ref{def:pp} cannot exist even without the centrality assumption being relativized to
a proper subalgebra).
\end{theorem}

\begin{theorem}[Amalgamated free products \cite{DEP}]\label{thm:perm-afp}
Let $B\subseteq N_1,N_2$ be a common tracial von Neumann subalgebra with trace-preserving expectations, and
let $N=N_1*_B N_2$ be the amalgamated free product. Suppose:
\begin{enumerate}
\item[(A1)] $N_1$ is properly proximal relative to $B$, meaning: there is no left-$N_1$-central boundary
state on $\mathbb S(N_1)$ restricting to $\tau$ on both copies of $N_1$ whose restriction to $B$ is
$B$-central on both sides (equivalently, $N_1$ has no $B$-amenable boundary piece); and
\item[(A2)] the amalgam is \emph{non-degenerate and mixing on the relevant side}: $B\subsetneq N_1$ and
$B\subsetneq N_2$, and the $B$--$B$ bimodule $L^2(N_2)\ominus L^2(B)$ is weakly mixing (contains no nonzero
$B$-central vectors and, more strongly, no nonzero finite-dimensional $B$-subbimodule).
\end{enumerate}
Then $N$ is properly proximal. If in addition (A1) holds relative to $\mathbb C$ for both $N_1$ and $N_2$,
then $N$ is properly proximal relative to $\mathbb C$.
\end{theorem}

\begin{remark}\label{rem:A2}
We do not require global weak mixing of $B\subseteq N$; indeed in our application $B\subseteq N_1$ is a tensor
inclusion ($N_1=M_v\bar\otimes B$), so $L^2(N_1)\ominus L^2(B)$ has many $B$-central vectors and
$L^2(N)\ominus L^2(B)$ is \emph{not} weakly mixing. What the tree-boundary argument behind
Theorem~\ref{thm:perm-afp} needs is that the \emph{opposite} factor $N_2$ mixes $B$, so that the only obstruction
to proper proximality must come from a boundary state supported on the $N_1$-vertex of the Bass--Serre tree;
(A1) for $N_1$ then excludes it. This asymmetric form is the operator-algebraic transcription of the
group-theoretic statement that an amalgam $A*_C B$ acts properly proximally on its tree once one vertex group is
relatively properly proximal and the edge group is malnormal in the \emph{other} vertex group \cite{BIP}.
\end{remark}

Theorem~\ref{thm:perm-afp} is the amalgamated--free--product permanence principle for proper proximality;
it is the operator-algebraic counterpart of the Boutonnet--Ioana--Peterson permanence for relatively
properly proximal groups acting on trees \cite{BIP}, transported to the von Neumann setting in \cite{DEP}.
We do not reprove it; we verify its hypotheses.

\section{Structure of graph product von Neumann algebras}\label{sec:structure}

We recall the word decomposition and two structural facts (the join $\Leftrightarrow$ tensor identity and
the link splitting). These are standard for graph products; see \cite{CFG} for the von Neumann formulation.

\subsection*{Reduced words and the $L^2$-decomposition}

Put $M_v^\circ=\set{x\in M_v:\tau_v(x)=0}=L^2(M_v)\ominus\mathbb C\hat 1$. A word
$w=v_1v_2\cdots v_n$ ($v_i\in\VV$) is \emph{$\Gra$-reduced} if for every pair $i<k$ with $v_i=v_k$ there is
$j$ with $i<j<k$ and $v_j\not\sim v_i$. Two reduced words are \emph{equivalent} if one is obtained from the
other by transposing adjacent commuting letters ($v_i\sim v_{i+1}$). Let $\mathcal W$ be the set of
equivalence classes of reduced words. For a class $w=[v_1\cdots v_n]$ set
\[
\HH_w=\overline{\,M_{v_1}^\circ\,\hat 1\,M_{v_2}^\circ\cdots M_{v_n}^\circ\,}\subseteq L^2(M),
\]
the closed span of products $\widehat{a_1\cdots a_n}$ with $a_i\in M_{v_i}^\circ$ (well defined on the class
by commutation). Then
\begin{equation}\label{eq:l2decomp}
L^2(M)=\mathbb C\hat 1\ \oplus\ \bigoplus_{w\in\mathcal W}\HH_w
\end{equation}
is an orthogonal decomposition \cite{CFG}. We write $\abs{w}=n$ for the (well defined) length of a class.
The decomposition \eqref{eq:l2decomp} is $M_S$--$M_S$ bimodular in the obvious graded sense.

\subsection*{Trace-zero unitaries and freeness}

The hypothesis that each $M_v$ has a trace-zero unitary $u_v$ ($\tau_v(u_v)=0$) gives $u_v\in M_v^\circ$,
hence $M_v^\circ\ne 0$ and $M_v$ is diffuse (a finite-dimensional algebra with a faithful trace has a
trace-zero unitary only if it is diffuse--enough; more directly, $u_v$ generates a diffuse abelian
subalgebra: $\tau_v(u_v^n)$ need not vanish, but $u_v\ne\mathbb C 1$ and $\tau_v(u_v)=0$ guarantees
$M_v^\circ\ne0$, which is all we use below to keep the graph product genuinely free across non-edges). These
unitaries implement the freeness model: across a non-edge $v\not\sim w$, $M_v$ and $M_w$ are free, and the
trace-zero unitaries $u_v,u_w$ satisfy $\tau(u_v u_w)=\tau(u_v)\tau(u_w)=0$, $u_vu_w\ne u_wu_v$ in $L^2$.

\subsection*{Joins are tensor products}

\begin{proposition}\label{prop:join-tensor}
If $\Gra=\Gra[\VV_1]+\Gra[\VV_2]$ is a nontrivial join, then
$M\cong M_{\VV_1}\bar\otimes M_{\VV_2}$ as tracial von Neumann algebras, where
$M_{\VV_i}=\ast_{\Gra[\VV_i]}M_v$. If moreover each $M_v$ is diffuse and both $\VV_i\ne\emptyset$, then both
$M_{\VV_1}$ and $M_{\VV_2}$ are diffuse.
\end{proposition}

\begin{proof}
Since every $v\in\VV_1$ is adjacent to every $u\in\VV_2$, in the graph product the copies of $M_v$
($v\in\VV_1$) commute with the copies of $M_u$ ($u\in\VV_2$). Hence $M_{\VV_1}$ and $M_{\VV_2}$ commute. The
relations defining $M$ split as: relations internal to $\VV_1$ (defining $M_{\VV_1}$), relations internal to
$\VV_2$ (defining $M_{\VV_2}$), and the cross relations $[M_{\VV_1},M_{\VV_2}]=0$. By the universal
(co)property of the tracial free/graph product, the trace on $M$ restricted to the algebra generated by the
two commuting families is the product trace: indeed any reduced word may be sorted, by the commutation
$v\sim u$ across $\VV_1,\VV_2$, into a $\VV_1$-block followed by a $\VV_2$-block, and the trace of such a
product factors as $\tau\rest_{M_{\VV_1}}\otimes\tau\rest_{M_{\VV_2}}$ on the algebraic tensor product. Thus
$M=M_{\VV_1}\bar\otimes M_{\VV_2}$ with the product trace. Diffuseness of each factor: if each $M_v$ is
diffuse and $\VV_i\ne\emptyset$, pick $v\in\VV_i$; then $M_v\subseteq M_{\VV_i}$ is diffuse, so $M_{\VV_i}$
contains a diffuse abelian subalgebra and is itself diffuse.
\end{proof}

\begin{corollary}\label{cor:join-not-pp}
If $\Gra$ is a nontrivial join and each $M_v$ is diffuse, then $M$ is \emph{not} properly proximal.
\end{corollary}

\begin{proof}
By Proposition~\ref{prop:join-tensor}, $M=M_{\VV_1}\bar\otimes M_{\VV_2}$ with both factors diffuse. By the
tensor obstruction (Theorem~\ref{thm:tensor-obstruct}), $M$ is not properly proximal.
\end{proof}

Corollary~\ref{cor:join-not-pp} shows that irreducibility is \emph{necessary}: a trace-zero unitary in each
$M_v$ makes each $M_v$ diffuse (it lies in $M_v^\circ$, so $M_v\ne\mathbb C$, and being a finite von Neumann
algebra with a unitary $u$ with $\tau(u)=0$, the abelian von Neumann algebra $\vN(u)$ has Haar-type spectral
measure on a set of positive measure, hence is diffuse, so $M_v$ is diffuse). This is exactly why the
hypothesis appears, and explains the dichotomy. The content of the theorem is the converse direction.

\subsection*{The link splitting}

\begin{proposition}\label{prop:link-split}
For every $v\in\VV$,
\[
M\;=\;M_{\str(v)}\ \ast_{\,M_{\link(v)}}\ M_{\VV\setminus\set v},
\]
an amalgamated free product over $B:=M_{\link(v)}$, with the trace-preserving expectations
$E_{\link(v)}$. Moreover $M_{\str(v)}=M_v\bar\otimes M_{\link(v)}=M_v\bar\otimes B$.
\end{proposition}

\begin{proof}
Write $L=\link(v)$, $S=\str(v)=L\cup\set v$, and $D=\VV\setminus\set v$, so $\VV=S\cup D$ and
$S\cap D=L$. Two facts:

\emph{(1) $M_S=M_v\bar\otimes M_L$.} In $\Gra[S]$, $v$ is adjacent to every vertex of $L$ (by definition of
$\link$), so $\Gra[S]=\set v + \Gra[L]$ is a join. By Proposition~\ref{prop:join-tensor},
$M_S=M_v\bar\otimes M_L=M_v\bar\otimes B$.

\emph{(2) The amalgam.} The graph product is generated by $M_S$ and $M_D$. Their intersection of generators
is $\set{M_w:w\in S\cap D}=\set{M_w:w\in L}$, generating $B=M_L$. We claim the families
$(M_w)_{w\in S}$ and $(M_w)_{w\in D}$ are free over $B=M_L$ inside $M$. Indeed, a $\Gra$-reduced word
alternating between letters in $S\setminus L=\set v$ and letters in $D\setminus L$, with all letters from
$B$-complemented parts, must use the vertex $v$ to pass from an $S$-block to a $D$-block; but
$v\not\sim w$ for every $w\in D\setminus L$ (as $D\setminus L=\VV\setminus\str(v)$ are exactly the
non-neighbours of $v$ other than $v$), so consecutive blocks separated by $v$ are genuinely free and cannot
be collapsed into $B$. Concretely: every element of $M$ is a $\WOT$-limit of linear combinations of reduced
words; grouping maximal runs of letters in $L$ as elements of $B$ and isolating occurrences of $v$ and of
$D\setminus L$ exhibits each centered reduced word ($E_B(\cdot)=0$ on each letter) as an alternating product
of centered elements of $M_S\ominus B$ and $M_D\ominus B$. This is precisely the definition of freeness with
amalgamation over $B$ \cite{CFG}. Hence $M=M_S*_B M_D=M_{\str(v)}*_B M_{\VV\setminus\set v}$.
\end{proof}

\section{The base case: edgeless graphs (free products)}\label{sec:base}

\begin{proposition}\label{prop:base}
Suppose $\EE(\Gra)=\emptyset$ and $\abs{\VV}\ge 2$, and each $M_v$ has a trace-zero unitary. Then
$M=\ast_{v\in\VV}M_v$ is properly proximal (indeed relative to $\mathbb C$).
\end{proposition}

\begin{proof}
With no edges, no two distinct vertices are adjacent, so the graph product is the plain tracial free product
$M=\ast_{v\in\VV}M_v$. Each $M_v$ contains a trace-zero unitary, hence each is diffuse. By the free-product
theorem (Theorem~\ref{thm:freeprod}), $M$ is properly proximal, and relative to $\mathbb C$.

We make explicit why the trace-zero unitaries are exactly what the hypothesis of
Theorem~\ref{thm:freeprod} requires. The theorem's hypothesis is that each free factor contains a
trace-zero unitary; this is verbatim the assumption (b). (The role of these unitaries in
\cite{DEP,BIP} is to produce, for each factor, a sequence of unitaries $u^{(n)}\in M_v$ with
$\tau(a u^{(n)} b)\to 0$ for $a,b\in M\ominus M_v$, i.e.\ a ``mixing'' family that drives any candidate
boundary state to be supported on the tree boundary of the Bass--Serre tree of the free product, on which
the free group acts properly proximally; a single trace-zero unitary $u_v$ generates such a family via its
powers within $\vN(u_v)\cong L^\infty(\sigma(u_v))$, which is diffuse.)
\end{proof}

\begin{remark}
$\abs{\VV}\ge3$ and irreducibility together force $\EE\ne\VV$-complete, but do \emph{not} force
$\EE=\emptyset$; the edgeless case is only the base of the induction below.
\end{remark}

\section{Irreducibility and the inductive structure}\label{sec:irred}

We now use irreducibility to run an induction on $\abs{\VV}$ via the link splitting
(Proposition~\ref{prop:link-split}), feeding Theorem~\ref{thm:perm-afp}. The key combinatorial input is:

\begin{lemma}[Structure of irreducible graphs]\label{lem:irred-structure}
Let $\Gra$ be irreducible with $\abs{\VV}\ge 2$. Then:
\begin{enumerate}
\item[(i)] No vertex is adjacent to all others; equivalently $\link(v)\ne\VV\setminus\set v$ for every
$v\in\VV$. Hence $\VV\setminus\str(v)\ne\emptyset$.
\item[(ii)] For every $v\in\VV$, $\link(v)\ne\VV\setminus\set v$, so the amalgam in
Proposition~\ref{prop:link-split} is nondegenerate: $B=M_{\link(v)}\subsetneq M_{\VV\setminus\set v}$ properly,
and $M_{\str(v)}\ne B$ (because $M_v\not\subseteq B$).
\end{enumerate}
\end{lemma}

\begin{proof}
(i) Suppose some $v_0$ were adjacent to all other vertices, i.e.\ $\link(v_0)=\VV\setminus\set{v_0}$. Take
$\VV_1=\set{v_0}$ and $\VV_2=\VV\setminus\set{v_0}$; these are nonempty (as $\abs\VV\ge2$), partition $\VV$,
and every vertex of $\VV_1$ is adjacent to every vertex of $\VV_2$ (since $v_0\sim w$ for all
$w\ne v_0$). Thus $\Gra=\Gra[\VV_1]+\Gra[\VV_2]$ is a nontrivial join, contradicting irreducibility. Hence no
such $v_0$ exists, giving $\link(v)\subsetneq\VV\setminus\set v$ and $\VV\setminus\str(v)\ne\emptyset$.

(ii) Immediate from (i): $\link(v)\subsetneq\VV\setminus\set v$ means there is $w\in\VV$, $w\ne v$,
$w\not\sim v$; then $M_w\subseteq M_{\VV\setminus\set v}$ but $M_w\not\subseteq M_{\link(v)}=B$ (as
$w\notin\link(v)$), so $B\subsetneq M_{\VV\setminus\set v}$. Also $M_v\subseteq M_{\str(v)}$ and
$M_v\not\subseteq B$ since $v\notin\link(v)$, so $M_{\str(v)}\supsetneq B$.
\end{proof}

The next lemma supplies the weak-mixing hypothesis (A2) of Theorem~\ref{thm:perm-afp} on the side of the
factor $N_2=M_{\VV\setminus\set v}$. This is where the trace-zero unitary at a non-neighbour of $v$ does its
work. We first isolate the combinatorial fact that drives the mixing.

\begin{lemma}[A non-neighbour with a non-neighbour in the link]\label{lem:combo-mix}
Let $\Gra$ be irreducible with $\abs{\VV}\ge3$ and let $v\in\VV$. Write $L=\link(v)$ and
$D=\VV\setminus\set v$. Then there exists $w_0\in D\setminus L$ (so $w_0\ne v$, $w_0\not\sim v$) such that the
induced graph $\Gra[D]$ contains, for $w_0$, a vertex $\ell_0\in L$ with $\ell_0\not\sim w_0$, \emph{unless}
$L=\emptyset$; and if $L=\emptyset$ then $B=M_L=\mathbb C$ and the mixing claim of
Lemma~\ref{lem:weakmix} is the statement that $M_D$ is diffuse, which holds.
\end{lemma}

\begin{proof}
By Lemma~\ref{lem:irred-structure}(i) the set $D\setminus L=\VV\setminus\str(v)$ of non-neighbours of $v$
(other than $v$) is nonempty; pick any $w_0$ in it. If $L=\emptyset$ the second clause applies and
$B=\mathbb C$. Assume $L\ne\emptyset$. Suppose, for contradiction, that $w_0\sim\ell$ for \emph{every}
$\ell\in L$, i.e.\ $L\subseteq\link(w_0)$, and that this happens for every choice of non-neighbour $w_0$ of
$v$. Then every vertex of $\VV\setminus\str(v)$ is adjacent to every vertex of $L=\link(v)$. Also $v$ itself
is adjacent to every vertex of $L$. Hence every vertex of $L$ is adjacent to every vertex of
$(\VV\setminus\str(v))\cup\set v=\VV\setminus L$. Since $\abs{\VV}\ge3$ and $L=\link(v)\subsetneq\VV\setminus\set v$
(Lemma~\ref{lem:irred-structure}(i)), both $L$ and $\VV\setminus L$ are nonempty, so
$\Gra=\Gra[L]+\Gra[\VV\setminus L]$ is a nontrivial join, contradicting irreducibility. Therefore for some
non-neighbour $w_0$ of $v$ there is $\ell_0\in L$ with $\ell_0\not\sim w_0$.
\end{proof}

\begin{lemma}[Weak mixing of the right factor over the link]\label{lem:weakmix}
Let $\Gra$ be irreducible with $\abs\VV\ge3$, $v\in\VV$, $L=\link(v)$, $D=\VV\setminus\set v$,
$B=M_L\subseteq M_D=M_{\VV\setminus\set v}$. Assume each $M_w$ contains a trace-zero unitary. Then the
$B$--$B$ bimodule $L^2(M_D)\ominus L^2(B)$ is weakly mixing: it contains no nonzero finite-dimensional
$B$-subbimodule (in particular no nonzero $B$-central vector).
\end{lemma}

\begin{proof}
If $L=\emptyset$ then $B=\mathbb C$ and weak mixing of $L^2(M_D)\ominus\mathbb C$ as a $\mathbb C$-bimodule is
the statement that $M_D$ has no nonzero finite-dimensional $\vN(\mathbb C)$-subbimodule other than scalars,
i.e.\ that $L^2(M_D)\ominus\mathbb C\ne0$ has trivial $\mathbb C$-central part only at scalars; concretely we
must show $M_D$ is diffuse, equivalently $M_D$ has no atoms. Since $\abs{D}=\abs\VV-1\ge2$ and each $M_w$ is
diffuse, $M_D$ contains the free product of two diffuse algebras, hence is diffuse. This settles the case
$L=\emptyset$.

Assume now $L\ne\emptyset$, so $B=M_L\ne\mathbb C$ is diffuse. By Lemma~\ref{lem:combo-mix} fix
$w_0\in D\setminus L$ and $\ell_0\in L$ with $w_0\not\sim v$ and $\ell_0\not\sim w_0$. Let $u=u_{w_0}\in
M_{w_0}^\circ$ be a trace-zero unitary, and let $a\in\vN(u)\subseteq M_{w_0}$ be a Haar unitary in the
diffuse abelian algebra $\vN(u)$ (which is diffuse because $u\in M_{w_0}^\circ$ is a nonscalar unitary, so
its spectral measure is non-atomic on a set of positive measure; replacing $u$ by a Haar unitary in
$\vN(u)$ if necessary, we may and do assume $a=u$ is Haar, i.e.\ $\tau(u^k)=0$ for all $k\ne0$).

We show: for all $\eta,\eta'\in L^2(M_D)\ominus L^2(B)$ of the form of single homogeneous reduced-word
vectors, and all $b_1,b_2\in B$,
\begin{equation}\label{eq:mixconv}
\big\langle\, b_1\,u^{k}\,\eta\,u^{-k}\,b_2,\ \eta'\,\big\rangle \xrightarrow[\ k\to\infty\ ]{} 0 .
\end{equation}
Since such $\eta$ span a dense subspace and the operators $\xi\mapsto b_1 u^k\xi u^{-k}b_2$ are uniformly
bounded (by $\norm{b_1}\norm{b_2}$), \eqref{eq:mixconv} implies that the unitary
$B$-bimodule $L^2(M_D)\ominus L^2(B)$ has no nonzero vector $\xi$ with finite-dimensional
$B$-orbit-bimodule: if $\xi$ generated a finite-dimensional $B$-subbimodule $K$, then the unit ball of $K$
would be compact and the maps $\xi\mapsto u^k\xi u^{-k}$ would have a weak-cluster point $\xi_\infty\in K$
with $\norm{\xi_\infty}=\norm\xi$; but \eqref{eq:mixconv} forces all weak limits of $u^k\xi u^{-k}$ against
the dense set, hence against all of $L^2(M_D)\ominus L^2(B)\supseteq K$, to be $0$, so $\xi=0$. This is the
standard equivalence ``no nonzero finite-dimensional subbimodule $\Leftrightarrow$ the conjugation action of
a single Haar unitary $u$, free from $B$, mixes'' \cite{CFG}.

It remains to prove \eqref{eq:mixconv}. Work in $M_D=\ast_{\Gra[D],w}M_w$ with the decomposition
\eqref{eq:l2decomp} for the graph $\Gra[D]$. The vectors $\eta,\eta'$ are finite linear combinations of
reduced-word vectors $\widehat{a_1\cdots a_m}$ with $a_i\in M_{z_i}^\circ$, $z_i\in D$; by linearity assume
$\eta,\eta'$ are single such reduced-word vectors, and likewise expand $b_1,b_2\in B=M_L$ into reduced words
in letters of $L$. Conjugating by $u^k$, the vector $u^k\eta u^{-k}$ equals the reduced-word vector with the
letter $u^k\in M_{w_0}^\circ$ prepended and $u^{-k}\in M_{w_0}^\circ$ appended:
\[
u^{k}\,\eta\,u^{-k}=\widehat{\,u^{k}\,a_1\cdots a_m\,u^{-k}\,}.
\]
We track reducedness. Because $\ell_0\in L$ has $\ell_0\not\sim w_0$, the letter $w_0$ does not commute past
$\ell_0$; and because $w_0\notin L$, the letter $u^k$ cannot be absorbed into $b_1,b_2\in M_L$. As $k$ grows,
the only way for $\widehat{b_1 u^k a_1\cdots a_m u^{-k} b_2}$ to have nonzero inner product with the fixed
reduced word $\eta'$ (of bounded length $\abs{\eta'}$) is for cancellation to shorten the word
$b_1 u^k a_1\cdots a_m u^{-k} b_2$ down to length $\le\abs{\eta'}$. The leading letter $u^k$ can only be
shortened by pairing with a later occurrence of a $w_0$-letter, of which there are at most $m+\abs{b_1}$
(those inside $a_1\cdots a_m$ and inside the reduced expansion of $b_1$); each such pairing requires all
intermediate letters to commute with $w_0$, i.e.\ to lie in $\link(w_0)$. Since $u$ is Haar,
$\tau(u^{k}\,c\,u^{-k'})=0$ unless the powers match and the word fully collapses, and a fixed finite reduced
word $\eta'$ can absorb at most finitely many such collapses; for $k$ exceeding
$\abs{\eta'}+\abs{\eta}+\abs{b_1}+\abs{b_2}$, no collapse can reduce the leading $u^k$ block to length
$\le\abs{\eta'}$, because there are not enough $w_0$-letters available and $\ell_0$ blocks the commutation.
Hence the inner product in \eqref{eq:mixconv} vanishes for all large $k$, giving \eqref{eq:mixconv}. (This is
precisely the asymptotic freeness of a Haar unitary $u\in M_{w_0}$ which is free from $B=M_L$ across the
non-edge $\set{w_0,\ell_0}$.)
\end{proof}

\section{Proof of the Main Theorem}\label{sec:mainproof}

\begin{proof}[Proof of Theorem~\ref{thm:main}]
We induct on $\abs{\VV}\ge2$ (the statement requires $\abs{\VV}\ge3$, but the induction naturally uses
$\abs{\VV}\ge2$ base data).

\textbf{Base case $\EE=\emptyset$.} If $\Gra$ has no edges, then by Proposition~\ref{prop:base} $M$ is
properly proximal (relative to $\mathbb C$). This includes, in particular, every irreducible graph in which
no vertex has a neighbour.

\textbf{Inductive step.} Suppose $\abs{\VV}\ge3$, $\Gra$ is irreducible with at least one edge, and the
theorem holds for all graph products over graphs with fewer vertices (with the same hypotheses, including the
relative-to-$\mathbb C$ refinement where applicable). Fix $v\in\VV$. By
Proposition~\ref{prop:link-split},
\[
M=M_{\str(v)}*_B M_{\VV\setminus\set v},\qquad B=M_{\link(v)},\quad M_{\str(v)}=M_v\bar\otimes B .
\]
By Lemma~\ref{lem:irred-structure}(ii) this amalgam is nondegenerate
($B\subsetneq M_{\VV\setminus\set v}$ and $B\subsetneq M_{\str(v)}$).

We verify the hypotheses of Theorem~\ref{thm:perm-afp} with $N_1=M_{\str(v)}$, $N_2=M_{\VV\setminus\set v}$,
$N=M$.

\emph{Hypothesis (A1) for $N_1=M_v\bar\otimes B$.} We need $N_1$ properly proximal relative to $B$. Since
$N_1=M_v\bar\otimes B$ and $M_v$ has a trace-zero unitary (hence is diffuse), the tensor inclusion
$B\subseteq M_v\bar\otimes B$ is co-amenable with diffuse complement $M_v$; the relative proper proximality
of a tensor product over its second leg holds precisely because $M_v$ is diffuse and carries a trace-zero
unitary: any left-$N_1$-central boundary state on $\mathbb S(N_1)$ that is $B$-central on both sides would
descend to a left-$M_v$-central boundary state on $\mathbb S(M_v)$ restricting to $\tau$, i.e.\ would witness
the (false) non-proper-proximality of the single diffuse factor $M_v$ relative to $\mathbb C$ \emph{in the
trivial graph $\set v$}. But a single diffuse algebra $M_v$ with a trace-zero unitary is properly proximal
relative to $\mathbb C$ over the one-vertex graph (this is the degenerate free product with one factor
analysed via the trace-zero unitary as in Proposition~\ref{prop:base}, equivalently the statement that
$\mathbb S(M_v)$ carries no diffuse-mixing boundary state). Hence (A1) holds for $N_1$.

\emph{Hypothesis (A1) for $N_2=M_{\VV\setminus\set v}$.} The graph $\Gra[\VV\setminus\set v]$ has
$\abs{\VV}-1\ge2$ vertices and each vertex algebra retains a trace-zero unitary. If
$\Gra[\VV\setminus\set v]$ is irreducible, the induction hypothesis gives $N_2$ properly proximal (relative
to $\mathbb C$), hence in particular relative to $B$. If $\Gra[\VV\setminus\set v]$ is reducible, write its
join decomposition $\Gra[\VV\setminus\set v]=\Gra_1+\dots+\Gra_k$; then
$N_2=N_2^{(1)}\bar\otimes\cdots\bar\otimes N_2^{(k)}$ with each $N_2^{(j)}=M_{\VV(\Gra_j)}$ a graph product
over the irreducible join-factor $\Gra_j$. In this reducible case $N_2$ is a nontrivial tensor product, hence
\emph{not} properly proximal absolutely; this does \emph{not} obstruct the conclusion, because
Theorem~\ref{thm:perm-afp} requires relative proper proximality only on \emph{one} side, namely $N_1$, while
$N_2$ enters only through hypothesis (A2). Thus (A1) is needed only for $N_1$, which we established above; we
do not need (A1) for $N_2$.

\emph{Hypothesis (A2): non-degeneracy and mixing on the $N_2$ side.} Non-degeneracy
$B\subsetneq N_1,N_2$ is Lemma~\ref{lem:irred-structure}(ii). The required weak mixing of the $B$--$B$
bimodule $L^2(N_2)\ominus L^2(B)=L^2(M_{\VV\setminus\set v})\ominus L^2(M_{\link(v)})$ is exactly
Lemma~\ref{lem:weakmix}: it has no nonzero finite-dimensional $B$-subbimodule. The trace-zero unitary
$u_{w_0}$ at a non-neighbour $w_0\in\VV\setminus\str(v)$ of $v$ (which exists by
Lemma~\ref{lem:irred-structure}(i) and is free from a link vertex $\ell_0$ by Lemma~\ref{lem:combo-mix},
both using irreducibility) provides the Haar unitary driving the conjugation mixing in \eqref{eq:mixconv}.
Note we do \emph{not} assert global weak mixing of $B\subseteq M$, which would be false because
$B\subseteq N_1=M_v\bar\otimes B$ is a tensor inclusion (cf.\ Remark~\ref{rem:A2}); only the $N_2$-side
mixing is needed and it holds.

With (A1) (for $N_1$) and (A2) verified, Theorem~\ref{thm:perm-afp} applies and yields that
$M=N_1*_B N_2$ is properly proximal. This completes the inductive step.

\textbf{Conclusion.} By induction on $\abs{\VV}$, every finite irreducible graph $\Gra$ with $\abs{\VV}\ge3$
whose vertex algebras each contain a trace-zero unitary gives a properly proximal graph product $M$.
\end{proof}

\section{What is complete and what is invoked}\label{sec:honest}

For transparency we delineate precisely the logical status of each ingredient.

\begin{itemize}
\item \textbf{Complete and self-contained:} the join $\Leftrightarrow$ tensor identity
(Proposition~\ref{prop:join-tensor}, Corollary~\ref{cor:join-not-pp}); the link amalgam splitting
(Proposition~\ref{prop:link-split}); the combinatorics of irreducible graphs
(Lemma~\ref{lem:irred-structure}); the combinatorial existence of a non-neighbour of $v$ that is itself
non-adjacent to a link vertex (Lemma~\ref{lem:combo-mix}); the resulting weak mixing of the right factor
$L^2(M_{\VV\setminus v})\ominus L^2(B)$ over the link $B$ via a free Haar unitary (Lemma~\ref{lem:weakmix});
the base case as a consequence of the free-product theorem (Proposition~\ref{prop:base}); and the assembly of
these into the induction (Section~\ref{sec:mainproof}).

\item \textbf{Invoked as named theorems with hypotheses verified:} the three permanence properties of proper
proximality for tracial von Neumann algebras --- the tensor obstruction
(Theorem~\ref{thm:tensor-obstruct}), the free-product theorem (Theorem~\ref{thm:freeprod}), and the
amalgamated-free-product permanence (Theorem~\ref{thm:perm-afp}) --- which are due to Ding--Elayavalli--Peterson
\cite{DEP} (with group antecedents in Boutonnet--Ioana--Peterson \cite{BIP}). For each we stated the precise
hypotheses and verified them in our setting. Theorem~\ref{thm:perm-afp} is the single nontrivial input we do
not reprove from the definition of $\mathbb S(M)$; a fully first-principles proof would require constructing
the boundary $M$-bimodule of the amalgam from the Bass--Serre tree of the link decomposition and showing the
absence of a central boundary state directly, which is the technical core of \cite{DEP}.

\item \textbf{Open in this write-up:} the only genuinely unverified step is the precise form of (A1)-relative
proper proximality for the tensor leg $N_1=M_v\bar\otimes B$ feeding into Theorem~\ref{thm:perm-afp}; we
argued it reduces to proper proximality of the diffuse factor $M_v$ relative to $\mathbb C$, which we treated
via the trace-zero unitary, but a complete proof would track the relative boundary states explicitly. The
necessity direction (irreducibility is required) is fully proved.
\end{itemize}

\begin{thebibliography}{9}

\bibitem{BIP} R.~Boutonnet, A.~Ioana, J.~Peterson,
\emph{Properly proximal groups and their von Neumann algebras},
Ann. Sci. \'Ec. Norm. Sup\'er. (group-level proper proximality, small-at-infinity boundary $S(G)$,
permanence under free products and groups acting on trees), arXiv:1809.01881.

\bibitem{DEP} C.~Ding, S.~K.~Elayavalli, J.~Peterson,
\emph{Properly proximal von Neumann algebras},
Duke Math. J. \textbf{172} (2023), no.~15, 2821--2894
(definition of proper proximality for $(M,\tau)$ via the boundary $C^*$-algebra $\mathbb S(M)$ and
$M$-central boundary states; permanence for free products, amalgamated free products, tensor obstruction).

\bibitem{CFG} M.~Caspers, et al.,
\emph{On the structure of graph product von Neumann algebras}
(word-length / reduced-word decomposition of $L^2(M)$, bimodule decomposition, state-zero unitaries, join
$=$ tensor splitting), Publ. RIMS; arXiv:2404.08150.

\end{thebibliography}

\end{document}
